\section*{अध्याय \ref{ch:b1:hotspots} --- हॉटस्पॉट और अंतःप्लेट ज्वालामुखी}

\begin{solution}{exo:b1:hotspots:1}
(क) $4205 + 5274 = \qty{9479}{m}$, लगभग \qty{9.5}{km}। ज्वालामुखी उस पानी से भारी है जिसे वह विस्थापित करता है: प्लेट उसके भार के नीचे मुड़ती है, अतः उसका आधार दूर के तल से गहरा बैठता है, और उसकी वास्तविक ऊँचाई अधिक है। (ख) $74000/0.1 = \num{740000}$ वर्ष किलाउआ की आपूर्ति पर; $74000/0.025 = \num{3.0e6}$ वर्ष शृंखला की औसत दर पर। (ग) $0.12\times19 = \qty{2.3}{km^3}$, माउना लोआ का लगभग $\num{3e-5}$: एक प्रतिशत का तीन हज़ारवाँ भाग।
\end{solution}

\begin{solution}{exo:b1:hotspots:2}
(क) $0.92\times\pi/180 = \qty{0.0161}{rad.Myr^{-1}}$। (ख) प्रति मिलियन वर्ष रेडियन गुणा किलोमीटर में, $\omega R\sin\Delta$ किलोमीटर प्रति मिलियन वर्ष में है, अर्थात् मिलीमीटर प्रति वर्ष: $0.016057\times6371 = 102.3$, अतः $v = 102.3\sin\Delta$, जो \ang{30} पर \qty{51.1}{mm.yr^{-1}}, \ang{60} पर \qty{88.6}{mm.yr^{-1}}, \ang{90} पर \qty{102.3}{mm.yr^{-1}} और \ang{108.0} पर \qty{97.3}{mm.yr^{-1}} देता है। (ग) चाल ध्रुव से परिभ्रमण-अक्ष से दूरी $R\sin\Delta$ के समानुपाती है, \omterm{def:b1:spherical-astronomy:spherical-triangle}{महावृत्त} पर ध्रुव से \ang{90} पर सबसे बड़ी; उसके आगे बिंदु फिर अक्ष के निकट आता है, विपरीत ध्रुव की ओर।
\end{solution}

\begin{solution}{exo:b1:hotspots:3}
(क) डेक्कन: $\num{e6}/\num{5e5} = \qty{2}{km}$। मध्य अटलांटिक: $\num{2e6}/\num{7e6} = \qty{0.29}{km}$, लगभग \qty{290}{m}। (ख) डेक्कन का माध्य उसके \qty{2.5}{km} खंड के निकट है: प्रांत एक मोटा, अविच्छिन्न पठार है। (ग) मध्य अटलांटिक का क्षेत्रफल उस पूरे क्षेत्र का है, चार महाद्वीपों पर, जिसके ऊपर उसके अवशेष मिलते हैं: लावा, डाइक और सिल, जो अब एक महासागर से अलग हैं और बहुत अपरदित हैं। चट्टान कुछ स्थानों पर संकेंद्रित है और अधिकांश क्षेत्र में अनुपस्थित है, अतः माध्य एक एकल ढेर के बारे में कुछ नहीं कहता।
\end{solution}

\begin{solution}{exo:b1:hotspots:4}
(क) $\rho_m\alpha\,\Delta T = 3300\times\num{3.2e-5}\times250 =
\qty{26.4}{kg.m^{-3}}$। (ख) स्तंभ-खंड में अतिरिक्त द्रव्यमान $(\rho_m - \rho_w)h = 2300\times690 = \qty{1.587e6}{kg.m^{-2}}$ है, जिसे हल्की परत को संतुलित करना होगा: $H = \num{1.587e6}/26.4 = \qty{6.0e4}{m}$, लगभग \qty{60}{km} गरम चट्टान।
\end{solution}

\begin{solution}{exo:b1:hotspots:5}
(क) A पर ($\lambda = 0$, $\varphi = 0$), $\vec e = (0, 1, 0)$ और
$\vec n = (0, 0, 1)$; $\sin\ang{270} = -1$ और $\cos\ang{270} = 0$ के साथ,
$\vec t_A = -\vec e = (0, -1, 0)$। B पर ($\lambda = \ang{90}$),
$\vec e = (-1, 0, 0)$, $\vec n = (0, 0, 1)$ और $\vec t_B = (1, 0, 0)$।
(ख) $\vec t_A\times\vec t_B = (0, 0, 1)$। (ग) $\vec p' = (0, 0, 1)$ और
$\vec r_A = (1, 0, 0)$ के साथ, $\vec p'\times\vec r_A = (0, 1, 0)$, और
$(\vec p'\times\vec r_A)\cdot\vec t_A = -1 < 0$: ध्रुव
$-\vec p' = (0, 0, -1)$ है, दक्षिण ध्रुव। भूमध्य-रेखा के साथ पश्चिम की ओर गति उत्तरी ध्रुव के ऊपर से देखने पर दक्षिणावर्त है, अतः दक्षिणी ध्रुव के ऊपर से देखने पर वामावर्त। (घ) ध्रुव से \ang{90} पर, $\omega = 50/6371 =
\qty{7.85e-3}{rad.Myr^{-1}} = \qty{0.45}{\degree\per\Myr}$।
\end{solution}

\begin{solution}{exo:b1:hotspots:6}
(क) $\sum x_i v_i = 90000 + 240000 + 342000 = \num{672000}$ और
$\sum x_i^2 = \num{7.0e7}$, अतः $\omega = \qty{0.0096}{rad.Myr^{-1}}$, या \qty{0.55}{\degree\per\Myr}। (ख) पूर्वानुमान 28.8, 48.0 और \qty{57.6}{mm.yr^{-1}} हैं; अवशेष $+1.2$, $0.0$ और $-0.6$ \unit{mm.yr^{-1}}। (ग) दृढ़ घूर्णन के लिए चाल अक्ष से दूरी के समानुपाती है, और अक्ष पर का बिंदु नहीं चलता: फ़िट करने के लिए कोई अंतःखंड नहीं है।
\end{solution}

\begin{solution}{exo:b1:hotspots:7}
(क) $v = \qty{50e-3}{m}/\qty{3.156e7}{s} = \qty{1.58e-9}{m.s^{-1}}$, अतः
$B = 2300\times\num{3e8}\times\num{1.58e-9} = \qty{1090}{kg.s^{-1}}$ और
$Q = (910/\num{3.2e-5})\times1090 = \qty{3.1e10}{W}$। (ख) हवाई प्रवाह का एक-तिहाई। (ग) $B$, $A\,v$ के साथ बढ़ता है: तेज़ प्लेट के नीचे का स्तंभ उसी उभार को दोगुनी तेज़ी से पुनः बनाना होगा, और उसका प्रवाह दोगुना बड़ा है।
\end{solution}

\begin{solution}{exo:b1:hotspots:8}
(क) $350/(16.1 - 10.3) = 350/5.8 = \qty{60}{km.Myr^{-1}}$, अर्थात् \qty{60}{mm.yr^{-1}}। (ख) युवा पथ, \qty{29 \pm 5}{mm.yr^{-1}}, \qty{26 \pm 5}{mm.yr^{-1}} वाली प्लेट के साथ अनिश्चितताओं के भीतर मेल खाता है; पुराना भाग दोगुने से अधिक तेज़ चला, $60 - 26 = 34$, प्लेट की चाल की लगभग सात मानक अनिश्चितताएँ। (ग) हाँ: N\,\ang{54}\,E दिगंश \ang{54} है, जो \ang{235} के एक अंश के भीतर उसके विपरीत है। ज्वालामुखी-क्रिया प्लेट के ऊपर उस दिशा के विपरीत चलती है जिसमें प्लेट \omterm{def:b1:hotspots:hotspot}{हॉटस्पॉट} के ऊपर चलती है। (घ) खिंचाव ने पुराने कैल्डेरों को अलग कर दिया, उनके बीच की दूरियाँ बढ़ाकर और इस प्रकार प्रत्यक्ष चाल; प्लेट स्वयं तेज़ नहीं चली।
\end{solution}

\begin{solution}{exo:b1:hotspots:9}
(क) $\num{9.2e10}/\num{e12} = \qty{0.092}{W.m^{-2}}$, \qty{92}{mW.m^{-2}}। (ख) प्रेक्षित अतिरिक्त का नौ से अठारह गुना, जो शायद शून्य भी हो। (ग) $L^2/\kappa =
\num{e10}/\num{e-6} = \qty{e16}{s}$, लगभग \qty{320}{Myr}; प्लेट \qty{1000}{km} को $1000/92 = \qty{11}{Myr}$ में पार करती है। जब प्लेट उभार के ऊपर है, नीचे से आई ऊष्मा को प्लेट के केवल एक छोटे भाग के आर-पार विसरित होने का समय मिला है: स्तंभ की अधिकांश ऊष्मा प्लेट के नीचे \omterm{def:b1:layered-earth:layers}{प्रावार} में रहती है और उसके साथ दूर ले जाई जाती है, और सतह उसमें से थोड़ी ही देखती है।
\end{solution}

\begin{solution}{exo:b1:hotspots:10}
(क) पिघले आयतन का दस गुना, \qty{e7}{km^3}, त्रिज्या $(3V/4\pi)^{1/3} = \qty{134}{km}$ का। (ख) $\num{e7}/100 = \qty{e5}{km^2}$, एक चकती जिसका व्यास $2\sqrt{\num{e5}/\pi} = \qty{357}{km}$ है: क्षेत्रफल में ट्रैप्स से अठारह गुना छोटी। शीर्ष बड़ा रहा होगा या कम पिघला होगा, या वह अधिक चौड़ा और पतला फैला होगा, या लावा उससे बहुत दूर तक बह गया। (ग) \qty{1}{km^3} प्रति वर्ष, किलाउआ की वर्तमान आपूर्ति का दस गुना। (घ) एक नए स्तंभ को श्यान \omterm{def:b1:layered-earth:layers}{प्रावार} के माध्यम से अपना रास्ता ऊपर धकेलने के लिए पर्याप्त गरम, हल्की चट्टान इकट्ठी करनी होगी, और एक बड़ा आयतन एक छोटे से कहीं अधिक तेज़ी से उठता है; उसके बाद आने वाली पुच्छ शीर्ष द्वारा खोले गए गरम, कमज़ोर मार्ग से ऊपर बहती है, और एक संकरी नली पर्याप्त है।
\end{solution}

\begin{solution}{exo:b1:hotspots:11}
(क) $\qty{40}{km.Myr^{-1}}\times(81 - 47)\,\unit{Myr} = \qty{1360}{km}$।
(ख) प्लेट $(-80\sin\ang{20}, 80\cos\ang{20}) = (-27.4, 75.2)$ पर चलती है और \omterm{def:b1:hotspots:hotspot}{हॉटस्पॉट} $(0, -40)$ पर: \omterm{def:b1:hotspots:hotspot}{हॉटस्पॉट} के सापेक्ष प्लेट
$(-27.4, 115.2)$ \unit{mm.yr^{-1}} पर चलती है, \qty{118.4}{mm.yr^{-1}} की चाल से
$\ang{360} - \arctan(27.4/115.2) = \ang{346.6}$ की ओर। (ग) पथ \qty{48}{\percent} अधिक तेज़ चाल और \ang{6.6} अधिक पूर्व की दिशा देता: खिसकता \omterm{def:b1:hotspots:hotspot}{हॉटस्पॉट} प्लेट की गति को स्वयं के सापेक्ष दर्ज करता है, गहरे \omterm{def:b1:layered-earth:layers}{प्रावार} के सापेक्ष नहीं।
\end{solution}

\begin{solution}{exo:b1:hotspots:12}
(क) $\qty{3e6}{J.m^{-3}.K^{-1}}\times\qty{1300}{K}\times\qty{e15}{m^3} =
\qty{3.9e24}{J}$। (ख) $\qty{3.16e13}{s}$ में फैलाने पर, \qty{1.2e11}{W}: लगभग वह ऊष्मा जो आज हवाई स्तंभ वाहित करता है (1.3 गुना), और पृथ्वी की \qty{47}{TW} का \qty{0.26}{\percent}। (ग) एक निचली सीमा: शीर्ष अपनी नौ-दसवीं चट्टान की अतिरिक्त ऊष्मा भी ऊपर लाता है जो नहीं पिघलती, और \omterm{def:b1:layered-earth:lithosphere}{स्थलमंडल} के आधार को गरम करता है, और गुप्त ऊष्मा छोड़ दी गई थी।
\end{solution}

\begin{solution}{pb:b1:hotspots:1}
\textbf{1.} हवाई पर ($\varphi = \ang{19.0}$, $\lambda = \ang{-155.2}$),
$\sin\ang{304} = -0.8290$ और $\cos\ang{304} = 0.5592$ के साथ:
\[
\begin{aligned}
  \vec r_H &= (-0.8583,\ -0.3966,\ 0.3256), &
  \vec e &= (0.4195,\ -0.9078,\ 0),\\
  \vec n &= (0.2955,\ 0.1366,\ 0.9455), &
  \vec t_H &= (-0.1825,\ 0.8289,\ 0.5287).
\end{aligned}
\]

\textbf{2.} लुईसविल पर ($\varphi = \ang{-53.6}$, $\lambda = \ang{-140.6}$),
$\sin\ang{316} = -0.6947$ और $\cos\ang{316} = 0.7193$ के साथ:
\[
\begin{aligned}
  \vec r_L &= (-0.4586,\ -0.3767,\ -0.8049), &
  \vec e &= (0.6347,\ -0.7727,\ 0),\\
  \vec n &= (-0.6220,\ -0.5109,\ 0.5934), &
  \vec t_L &= (-0.8883,\ 0.1693,\ 0.4269).
\end{aligned}
\]

\textbf{3.} $\lVert\vec t_H\rVert^2 = 0.0333 + 0.6871 + 0.2795 = 0.9999$ और
$\vec t_H\cdot\vec r_H = 0.1566 - 0.3287 + 0.1721 = 0.0000$, पूर्णांकन तक।

\textbf{4.} $\vec t_H\times\vec t_L = (0.2643, -0.3917, 0.7054)$, जिसकी लंबाई
$0.8491$ है। दोनों पथों के बीच का कोण $\arcsin0.8491 = \ang{58.1}$ है (उनका डॉट-गुणनफल, $0.528$, धनात्मक है, अतः कोण न्यून है)।

\textbf{5.} यह है
\[ \vec p' = (0.3113, -0.4613, 0.8308) \]
फिर इसके बाद $(\vec p'\times\vec r_H)\cdot\vec t_H =$ $-0.98 < 0$: ध्रुव
$\vec p = (-0.3113, 0.4613, -0.8308)$ है।

\textbf{6.} $\arcsin(-0.8308) = \ang{-56.18}$ और
$\operatorname{atan2}(0.4613, -0.3113) = \ang{124.01}$: ध्रुव \ang{56.2}\,S, \ang{124.0}\,E पर है, दक्षिणी महासागर में ऑस्ट्रेलिया के दक्षिण में।

\textbf{7.} $\vec p\cdot\vec r_H = 0.2672 - 0.1830 - 0.2705 = -0.1863$, अतः
$\Delta_H = \ang{100.7}$; $\vec p\cdot\vec r_L = 0.1428 - 0.1738 + 0.6687 =
0.6377$, अतः $\Delta_L = \ang{50.4}$।

\textbf{8.} सदिश $\vec p\times\vec r_L$, $\vec p$ और $\vec r_L$ के लंबवत है, और $\vec t_L$ भी है (जैसा $\vec p$ की रचना से, और क्योंकि पथ क्षैतिज है); दोनों एक ही रेखा के अनुदिश हैं, और उसकी लंबाई $\sin\Delta_L$ है। अतः $(\vec p\times\vec r_L)\cdot\vec t_L = \pm\sin\Delta_L$, और परिकलन देता है $+0.770 = \sin\ang{50.4}$: प्लेट लुईसविल पर $+\vec t_L$ के अनुदिश चलती है, पथ की दिशा में। चिह्न-नियम, जो हवाई पर लगाया गया, लुईसविल पर गलत अर्थ दे सकता था; उसने नहीं दिया, अतः दोनों दिशाएँ, \ang{304} और \ang{316}, ठीक दोहराई जाती हैं।

\textbf{9.} $\omega = 67/(6371\sin\ang{50.38}) = 67/4908 =
\qty{0.01365}{rad.Myr^{-1}}$, अर्थात् \qty{0.782}{\degree\per\Myr}।

\textbf{10.} $360/0.782 = \qty{460}{Myr}$।

\textbf{11.} \ang{56.18}\,S, \ang{124.01}\,E और \ang{63.78}\,S, \ang{112.48}\,E के बीच कोज्या-नियम \ang{9.5} देता है। दो-पथ दर, \qty{0.78}{\degree\per\Myr}, तेरह चालों पर फ़िट \qty{0.92}{\degree\per\Myr} से \qty{15}{\percent} कम है।

\textbf{12.} सोसाइटी \ang{73.6}, पिटकेर्न \ang{77.8}, कॉब \ang{134.7}।

\textbf{13.} $\omega R$ का मान:
\[ \omega R = 0.01365\times6371 = \qty{86.98}{mm.yr^{-1}} \]
गुणा से $\sin\Delta$:
सोसाइटी \qty{83.5}{mm.yr^{-1}};
पिटकेर्न \qty{85.0}{mm.yr^{-1}}, कॉब \qty{61.8}{mm.yr^{-1}}।

\textbf{14.} पिटकेर्न सहमत है: मापी गई से \qty{5}{mm.yr^{-1}} कम, अपनी \qty{15}{mm.yr^{-1}} के भीतर। सोसाइटी \qty{25.5}{mm.yr^{-1}} बहुत धीमी है, 2.6 अनिश्चितताएँ, और कॉब \qty{18.8}{mm.yr^{-1}} बहुत तेज़, छह से अधिक: कोई भी सहमत नहीं।

\textbf{15.} ध्रुव दोनों पथ-सदिशों के लंबवत है। यदि वे लगभग समांतर हों, तो उनका सदिश गुणनफल छोटा होता है, और किसी भी दिशा में कुछ अंशों की त्रुटि उसे बड़े कोण से घुमा देती है: ध्रुव खराब तय होता है।

\textbf{16.} $v = \qty{92e-3}{m}/\qty{3.156e7}{s} = \qty{2.92e-9}{m.s^{-1}}$
और $B = 2300\times\num{4.83e8}\times\num{2.92e-9} = \qty{3240}{kg.s^{-1}}$।

\textbf{17.} $Q = (910/\num{3.2e-5})\times3240 = \qty{9.2e10}{W}$,
\qty{0.20}{\percent} का \qty{47}{TW}।

\textbf{18.} कटक प्रति वर्ष प्रति किलोमीटर $21/67000 = \qty{3.1e-4}{km^3}$ बनाते हैं; $0.1/\num{3.1e-4} = \qty{320}{km}$ कटक।

\textbf{19.} $74000/0.1 = \num{740000}$ वर्ष।

\textbf{20.} $\num{2e6}/\num{81e6} = \qty{0.025}{km^3}$ प्रति वर्ष: किलाउआ को आज पूरे \omterm{def:b1:hotspots:hotspot}{हॉटस्पॉट} के औसत उत्पादन का चार गुना मिलता है, जो उसकी सबसे प्रबल अवस्था के निकट है।

\textbf{21.} $v_H = 86.98\sin\ang{100.73} = 86.98\times0.9825 =
\qty{85.5}{mm.yr^{-1}}$।

\textbf{22.} $92 - 85.5 = \qty{6.5}{mm.yr^{-1}}$, मापी गई चाल की 2.2 मानक अनिश्चितताएँ: दोनों \omterm{def:b1:hotspots:hotspot}{हॉटस्पॉट} एक स्थिर चौखटे पर एक \omterm{def:b1:plate-kinematics:rigid-plate}{दृढ़ प्लेट} की गति पर ठीक-ठीक सहमत नहीं हैं।

\textbf{23.} आयुएँ पूरी शृंखला पर औसत लेने पर \qty{9.1}{cm} और \qty{8.6}{cm} प्रति वर्ष देती हैं; मॉर्गन और मॉर्गन का \qty{9.2}{cm} प्रति वर्ष पहले के निकट है, और लुईसविल का पूर्वानुमान, \qty{8.5}{cm} प्रति वर्ष, दूसरे के निकट। चारों \qty{0.7}{cm} प्रति वर्ष के भीतर एक-दूसरे से हैं: पथ, आयुएँ और दिशाएँ लगभग \qty{8}{\percent} के भीतर एक ही कहानी कहते हैं।

\textbf{24.} एक दशमलव तक, लुईसविल पथ प्रशांत प्लेट को हवाई \omterm{def:b1:hotspots:hotspot}{हॉटस्पॉट} के ऊपर \textbf{\qty[text-series-to-math=true]{8.5}{cm.yr^{-1}}} की चाल देता है।
\end{solution}
